\documentclass[onecolumn]{autart}

%%%%%%%%%%%%%%%%%%%%%%%%%%%%%%%%%%%%%%%%%%%%%%%%%%%%%%%%%%%%%%%
% START CUSTOM INCLUDES & DEFINITIONS
%%%%%%%%%%%%%%%%%%%%%%%%%%%%%%%%%%%%%%%%%%%%%%%%%%%%%%%%%%%%%%%

%-------------------------------------------------------------
% 1. CORE MATH & FONT PACKAGES
%-------------------------------------------------------------
\usepackage{amsmath,amssymb,amsfonts}
\usepackage{mathtools}
\usepackage{mymath}
\usepackage{MnSymbol}

%-------------------------------------------------------------
% 2. REFERENCE & HYPERLINKING
%-------------------------------------------------------------
\usepackage[hidelinks]{hyperref}
\hypersetup{draft=false}

\usepackage[backend=bibtex, sorting=none, firstinits=true, abbreviate=true, url=false, isbn=false]{biblatex}
\addbibresource{master_bib_abbrev.bib}

%-------------------------------------------------------------
% 3. FIGURES, GRAPHS, & DRAWING (TikZ/PGF)
%-------------------------------------------------------------
\usepackage{tikz}
\usepackage{pgfplots}
\pgfplotsset{compat=newest}
\usetikzlibrary{plotmarks}
\usepgfplotslibrary{groupplots}

\usetikzlibrary{shapes, arrows, shapes.misc, arrows.meta, positioning, matrix, calc, fit, fadings, patterns}
\usetikzlibrary{calc,patterns,decorations.pathmorphing,decorations.markings,arrows, arrows.meta}
\usepackage{varwidth}
\usepackage[export]{adjustbox}
\usepackage{overpic}

%-------------------------------------------------------------
% 4. FLOATING OBJECTS, CAPTIONS, & LAYOUT
%-------------------------------------------------------------
\usepackage{caption}
\usepackage{subcaption}
\usepackage{placeins}

% That's because it likes placing large figures on their own page even though there's space for text
\renewcommand{\floatpagefraction}{0.7}
\renewcommand{\textfraction}{0.1}
\renewcommand{\topfraction}{0.9}
\renewcommand{\bottomfraction}{0.9}

%-------------------------------------------------------------
% 5. TABLES & ALIGNMENT
%-------------------------------------------------------------
\usepackage{arydshln}
\usepackage{tabto}
\usepackage{siunitx}
\usepackage{nth}

%-------------------------------------------------------------
% 6. ALGORITHMS & PSEUDOCODE
%-------------------------------------------------------------
\usepackage{algorithm}
\usepackage[noend]{algpseudocode}
\usepackage{xcolor}

\definecolor{commentgrey}{gray}{0.6}
\newcommand{\algcomment}[1]{\scriptsize\textcolor{commentgrey}{\Comment{#1}}}

%-------------------------------------------------------------
% 7. THEOREM-LIKE ENVIRONMENTS
%-------------------------------------------------------------
\newtheorem{assumption}{Assumption}
\newtheorem{remark}{Remark}

%-------------------------------------------------------------
% 8. UTILITY & CLASS FILE MODIFICATIONS
%-------------------------------------------------------------
\usepackage{todonotes}

% Remove preprint footer
\makeatletter
\def\ps@copyright{\let\@mkboth\@gobbletwo
\def\@oddhead{}
\let\@evenhead\@oddhead
\def\@oddfoot{\hfil\thepage\hfil}
\let\@evenfoot\@oddfoot}
\makeatother

% Patch to close NoHyper scope
\edef\endfrontmatter{%

\unexpanded\expandafter{\endfrontmatter}% the current code
\noexpand\endNoHyper
}

%-------------------------------------------------------------
% 9. CUSTOM MATH OPERATORS & COMMANDS
%-------------------------------------------------------------
\newcommand{\Argmin}{\ensuremath{\mathop{\mathrm{argmin}}}}
\newcommand{\Argmax}{\ensuremath{\mathop{\mathrm{argmax}}}}
\newcommand{\twonorm}[1]{\left\| #1 \right\|_2}
\newcommand{\set}[2]{\left\{ #1 \;\middle|\; #2 \right\}}
\newcommand{\inR}[3]{#1 \in \Rset^{#2 \times #3}}
\newcommand{\trans}[1]{#1^\top}

%%%%%%%%%%%%%%%%%%%%%%%%%%%%%%%%%%%%%%%%%%%%%%%%%%%%%%%%%%%%%%%
% END CUSTOM INCLUDES & DEFINITIONS
%%%%%%%%%%%%%%%%%%%%%%%%%%%%%%%%%%%%%%%%%%%%%%%%%%%%%%%%%%%%%%%

\pdfobjcompresslevel=0

\begin{document}
\begin{frontmatter}
\title{Fast-Forwarding Stalling in Dykstra's Algorithm\thanksref{footnoteinfo}}
\thanks[footnoteinfo]{Corresponding author I.~Kempf.}
\author[oxf]{Claudio Vestini}\ead{claudio.vestini@keble.ox.ac.uk}
and
\author[oxf]{Idris Kempf}\ead{idris.kempf@eng.ox.ac.uk}
\address[oxf]{Department of Engineering Science, University of Oxford, Parks Road, Oxford, OX1 3PJ, UK}


\begin{abstract}
Constrained quadratic programs and Euclidean projections are ubiquitous in engineering, arising in machine learning, estimation, control, and signal processing applications. Dykstra's algorithm is an iterative scheme for computing the Euclidean projection of an initial point onto the intersection of convex sets by successively projecting onto each set. Its simplicity and low computational cost make it well-suited for solving large-scale or real-time problems where traditional optimisation routines become computationally burdensome. Despite its strong convergence guarantees, Dykstra's algorithm is known to suffer from stalling---arbitrarily long intervals during which the primal iterates remain constant---rendering its runtime unpredictable and severely limiting its applicability in time-critical settings. Focusing on polyhedral constraint sets, we characterise the stalling phenomenon and derive a closed-form solution for the length of the stalling period once stalling is detected. This result enables a modified, stall-averse version of Dykstra's algorithm that fast-forwards the stalling period via a single, inexpensive update while preserving convergence guarantees. Numerical experiments demonstrate substantial improvements in convergence behaviour, establishing the proposed method as a practical enhancement for a broad class of projection-based algorithms.
\end{abstract}

\end{frontmatter}

\section{Introduction}

At the core of many problems in optimisation, statistics, and machine learning lies the task of identifying a point that satisfies multiple constraints simultaneously~\cite[Ch. 8]{BOYDCONVEX}. When these constraints correspond to closed convex sets in a Hilbert space, this task takes the form of the \textit{best approximation problem} (BAP). Formally, we denote as $\mathcal{X}$ a real Hilbert space equipped with a norm $\| \cdot \|$. We consider a finite family of closed convex subsets $\{\mathcal{H}_i\}_{i=0}^{n-1}\subseteq\mathcal{X}$ with a non-empty intersection $\mathcal{H} := \bigcap_{i=0}^{n-1} \mathcal{H}_i$. Given an arbitrary point $x^\circ \in \mathcal{X}$, the solution to the BAP is the unique point $x^\star \in \mathcal{H}$ closest to $x^\circ$ with respect to the norm~\cite{BAUSCHKEBOOK}. The BAP is then formulated as the convex optimisation problem $x^\star = \Argmin_{x\in \mathcal{H}} \frac{1}{2} \|x - x^\circ\|^2$. This formulation is widely adopted in applications such as restricted least squares and isotonic regression in statistics, as well as in signal and image processing, where convex sets encode prior knowledge, like non-negativity or sparsity, in control engineering, where system constraints are typically modeled using convex sets, and in machine learning, where solutions to learning problems must often adhere to structural constraints~\cite{BAUSCHKESWISS}.

In this paper, we assume that $\mathcal{X}$ is a Euclidean space, $X \subseteq \Rset^p$, equipped with the $\ell_2$ norm $\twonorm{x}\eqdef (x_1^2+\dots+x_p^2)^{1/2}$, so that the BAP takes the form of a \textit{constrained quadratic program} (CQP):
\begin{align}\label{eq:projection}
x^\star := \mathcal{P}_{\mathcal{H}}(x^\circ) := \argmin_{x\in \bigcap_{i=0}^{n-1}\mathcal{H}_i}\frac{1}{2}\twonorm{x-x^\circ}^2.
\end{align}
Here, we further assume that $\mathcal{H}$ is a polyhedral set, $\mathcal{H}\eqdef\set{x\in\Rset^p}{A x \leq c}$, and that each $\mathcal{H}_i$ is a half-space,
\begin{align}\label{eq:sets_i}
\mathcal{H}_i\coloneqq\set{x\in\Rset^p}{\trans{f_i}x \leq c_i}, \quad i=1,\dots,n,
\end{align}
where the normal vector $f_i$ of each plane $H_i\coloneqq\set{x\in\Rset^p}{\trans{f_i}x = c_i}$ satisfies $\twonorm{f_i}=1$. While in some cases the projection onto $\mathcal{H}$ is easy to compute, for arbitrary polyhedral sets the projection operator $\mathcal{P}_{\mathcal{H}}$ onto the intersection set is intractable and no closed-form solution to~\eqref{eq:projection} is known. In these cases, the solution to~\eqref{eq:projection} can be obtained using standard optimisation software packages. To solve~\eqref{eq:projection}, most solvers require introducing an additional variable $z\coloneqq Ax$ (e.g., \cite{osqp} or~\cite{ocpb:16}) that is projected onto the set $\set{z\in\Rset^n}{z \leq c}$, for which an explicit solution exists~\cite{BOYDCONVEX}. However, for large $n$, this variable augmentation can reduce the computational efficiency, particularly in high-dimensional or time-sensitive applications when the projection is part of an overarching algorithm~\cite{OPTIMNESTEROV,KEMPF20206542}.

% I would entirely remove this. The problem is that we're not talking about this in the remainder of the paper. If you'd re-run the fast-gradient using our modified Dyktra and show that it's faster you would need this paragraph.
% As a recent example, Diamond Light Source (DLS), the UK's national synchrotron, is expected to see an upgrade to its electron beam stabilisation control system. Existing control methods (modal decomposition~\cite{HEATH}) are no longer sufficient: this is due to the increased number of sensors (from 172 to 252) and actuators (from 173 to 396), as well as the introduction of two different types of corrector magnets. A devised model predictive control (MPC) implementation~\cite{MPCDLSII}, while promising, faces high computational demands for high-frequency (100 kHz) control. To accommodate input rate constraints, the Fast Gradient method was proposed in~\cite{KEMPF20206542}, where a CQP problem needs to be solved at each timestep as part of the main loop, which is to be run less than \SI{10}{\micro\second} for feasibility. Evaluating the solution to the CQP using standard solvers would be too expensive for this high-latency requirement. This difficulty necessitates iterative methods that construct the solution to the CQP in linear time using only the individual projection operators $\mathcal{P}_{\mathcal{H}_i}$. Amongst these, Dykstra's method offers a promising solution due to its simplicity and convergence guarantees.

As an alternative that can be faster in practice~\cite{KEMPF20206542}, \textit{Dykstra's projection algorithm}~\cite{DYKSTRA} is an iterative method designed to solve~\eqref{eq:projection}. It is an extension of the \textit{method of alternating projections} (MAP), which, given an initial point $x^\circ$, finds a point in the intersection of $\mathcal{H}_i$. Dykstra's algorithm modifies MAP by introducing auxiliary variables to achieve strong convergence to the specific projection of $x^\circ$ onto the intersection of $\mathcal{H}_i$ (or general convex sets). Qualitatively, Dykstra's auxiliary variables remember the direction vector of projection; these are added to the primary variables before subsequent projections and, as a result, the primary iterates of Dykstra's method asymptotically converge to~\eqref{eq:projection}. Although, like standard optimisation packages, Dykstra's method also introduces additional variables to solve~\eqref{eq:projection}, these appear only in vector additions rather than matrix-vector multiplications. This reduces the computational load and could make Dykstra's method a highly valuable algorithm to deploy in real-time, high-throughput applications for finding fast (approximate) solutions to CQPs. However, despite its strong theoretical guarantees, Dykstra's algorithm is known to suffer from a \emph{stalling} problem for certain classes of initial conditions~\cite{DYKSTRASTALLING}. In such cases, the iterates of Dykstra's method remain unchanged for a number of iterations. The duration of the stalling period is not known a priori, and, given the choice of starting point, can be made arbitrarily long. In general, the algorithm cannot be run for a fixed number of iterations with a guarantee that the output will be closer to the projection than the initial point, which can hinder deployment of Dykstra's method in practical applications.

In this paper, we address the stalling problem of Dykstra's method for polyhedral sets. Exploiting simplifications that arise in the polyhedral case, we derive the exact duration of the stalling period once stalling is detected (Theorem~\ref{thm:nstall}). This duration enables a modified version of Dykstra's algorithm (Algorithm~\ref{alg:fastforward}) that fast-forwards through the stalling period, resulting in improved convergence properties. Crucially, this modification preserves all convergence guarantees of the original algorithm (Corollary~\ref{cor:convergence}). This development represents a step towards practical deployment of Dykstra's algorithm in real-time settings, where its low-overhead projection steps could offer significant performance advantages.

The structure of the paper is as follows. Section~\ref{sec:background} reviews MAP, Dykstra's algorithm, and narrows the analysis to polyhedral sets, introducing the stalling phenomenon. Section~\ref{sec:main_result} proposes the fast-forward modification, and Section~\ref{sec:example} presents supporting numerical experiments, which demonstrate the effectiveness of the solution. Finally, Section~\ref{sec:conclusion} concludes with a discussion of potential accelerations, extensions, and applications. 

We use standard notation throughout: $\trans{A}$ denotes the transpose of vector or matrix $A$, and $\twonorm{A}$ its $\ell_{2}$-norm. The interior of a set $\mathcal{X}$ is denoted as $\text{int}\mathcal{X}$. Function composition is written as $f \circ g$, and the modulo and ceiling operators are denoted by $[\cdot]$ and $\lceil \cdot \rceil$, respectively.
\input{Latex/Current Version/Sections/2-Dykstra Background}
\section{Fast-Forwarding Dykstra's Algorithm}
\label{sec:main_result}

% In this section, we propose a modification to Dykstra's algorithm to resolve the stalling problem in the general case.

% \subsection{Stalling Formalisation}
Based on the discussion of stalling in~\cite{DYKSTRASTALLING}, we identify as a necessary condition for stalling of the iterates for successive iteration cycles over $n$ sets. This is formalised in Definition~\ref{def:stalling}.

\begin{defn}[Stalling]\label{def:stalling}
Given $m\geq n-1$, a stalling period is defined as those $i\in\lbrace 0,\dots,N_{stall}\rbrace$ with $N_{stall}\geq n$ for which $x_{m+i}=x_{m+i-n}$.
\end{defn}

During stalling, only the auxiliary variables $e_{m+i}$ change, which are modified by constant increments $x_m-x_{m+1}$. The stalling period therefore continues until $e_{m+i}$ becomes such that $x_{m+i}+e_{m-n+i}\in\mathcal{H}_{[m+i]}$. For polyhedral sets, the length of the stalling period can be pre-computed as soon as the algorithm encounters stalling, using the half-spaces that are active at this iteration. The stalling period continues until one of these half-spaces becomes inactive, at which point stalling terminates. This is formalised in Theorem~\ref{thm:nstall}.

\begin{thm}[Length of Stalling Period]\label{thm:nstall}
Suppose that stalling starts at iteration $m$ as in Definition~\ref{def:stalling} and denote the active half-spaces by $\mathcal{A}\coloneqq\set{i\in\lbrace 0,\dots,n-1\rbrace}{x_{m+i}+e_{m-n+i}\not\in\mathcal{H}_{[m+i]}}$. Then, the length of the stalling period is given by
\begin{align}\label{eq:nstall}
N_{stall}\coloneqq\min_{i\in\mathcal{S}} \left(i+\left\lceil k_{m-n+i} / \left(n(c_{[m+i]}-x_{m+i}^T f_{[m+i]})\right)\right\rceil\right),
\end{align}
where $k_{m-n+i}=e_{m-n+i}^T f_{[m+i]}$ and $\mathcal{S}\coloneqq\set{i\in \mathcal{A}}{x_{m-n+i}^T f_{[m+i]} - c_{[m+i]} < 0}$.
\end{thm}

\begin{pf}
Note that the existence of a finite number $N_{stall}$ is a consequence of the Boyle-Dykstra theorem. According to~\eqref{eq:dykstra:proj:poly}, the stalling period will terminate once an active half-space becomes inactive, i.e.\ $x_{m+i}+e_{m-n+i}\in\mathcal{H}_{[m+i]}$ for some $i\geq n$ (Definition~\ref{def:stalling}). According to~\eqref{eq:km}, the only half-spaces that can become inactive are those for which $x_{m+i}^T f_{[m+i]} - c_{[m+i]} < 0$, which is, by Definition~\ref{def:stalling}, a constant during stalling. The length of the stalling period can therefore be obtained from choosing the smallest possible integer $N_i$ for which
\begin{align*}
k_{m-n+i}+n N_i\left(x_{m+i}^T f_{[m+i]} - c_{[m+i]}\right) < 0,\qquad i\in\mathcal{A},
\end{align*}
so that $N_i=\left\lceil k_{m-n+i} / \left(n(c_{[m+i]}-x_{m+i}^T f_{[m+i]})\right)\right\rceil$. If $i_1<\dots<i_j$ are such that $N_{i_1}=\dots = N_{i_j}$, then according to~\eqref{eq:dykstra}, half-space $i_1$ is discarded first (and thus the minimiser of~\eqref{eq:nstall}).\qed
\end{pf}
Following the end of the stalling period, one half-space is discarded from the active half-space and $\mathcal{A}$ redefined, upon which the algorithm may again encounter a stalling condition. By Theorem~\eqref{thm:nstall}, the stalling period can be fast-forwarded by applying constant increments to $k_{m+i}$ for $i=0,\dots,N_{stall}$.

Theorem~\ref{thm:nstall} and Dykstra's algorithm can be combined as in Algorithm~\ref{alg:fastforward}. The algorithm proceeds identically to the standard Dykstra's algorithm until it detects stalling by verifying the stationarity condition $x_m = x_{m-n}$, as per Definition~\ref{def:stalling}. Once stalling is detected at iteration $m$, the algorithm executes its fast-forwarding phase. After performing the standard variable update for the current iteration, it computes the exact number of required stalling iterations, $N_{stall}$, by calling a function that implements Theorem~\ref{thm:nstall}. From this, it determines the number of full $n$-step cycles to skip, $N_{\text{cycles}} \coloneqq \lceil N_{stall} / n \rceil$. For each of the $n$ dual variables $e_j$, the algorithm then computes the constant increment $\Delta k_j = (\trans{f_j} x_{\text{stall}} - c_j)$ and applies a single, large update: $e_j \leftarrow e_j + N_{\text{cycles}} \cdot (\Delta k_j f_j)$. This ``jumps'' the dual variables forward to their state at iteration $m + N_{stall}$. Finally, the main iteration counter $m$ is advanced by $N_{stall}$, effectively skipping the entire stalling period. 

Since the proposed algorithm achieves the same iterates after escaping the stalling period (lines 13-18) as would be attained with standard Dykstra, the Boyle-Dykstra theorem~\cite{DYKSTRA} is retained, thus the iterates are guaranteed to converge to the optimal solution $x^\star$ asymptotically. This is formalised in Corollary~\ref{cor:convergence} below.

\begin{cor}[Convergence of Algorithm~\ref{alg:fastforward}]
\label{cor:convergence}
The sequence of iterates $\{x_m\}$ generated by the Stall-Averse Dykstra's Algorithm (Algorithm~\ref{alg:fastforward}) converges asymptotically to the optimal solution $x^\star = \mathcal{P}_{\mathcal{H}}(x^\circ)$, thus
\begin{align} \label{eq:boyle-dykstra}
    \lim_{m\rightarrow\infty}\|x_m-x^\star\|_2=0
\end{align}
\end{cor}

\begin{algorithm}
\caption{Stall-Averse Dykstra's Algorithm}
\label{alg:fastforward}
\begin{algorithmic}[1]
    \Function{DykstraFastForward}{$x^\circ, \{\mathcal{H}_i\}_{i=0}^{n-1}, N_{\max}, \epsilon_{\text{stall}}$}
        \State $x \leftarrow x^\circ$; $e_j \leftarrow 0$ \textbf{for} $j=0, \dots, n-1$; $X_{\text{hist}}[0] \leftarrow x^\circ$; $m \leftarrow 0$ \algcomment{Initialisation}
        \Statex\hrulefill
        \While{$m < N_{\max}$}
            \algcomment{Standard Dykstra's update}
            \State $e_{[m]} \leftarrow x + e_{[m-n]} - \mathcal{P}_{[m]}(x + e_{[m-n]})$
            \State $x \leftarrow x + e_{[m-n]} - e_{[m]}$ 
            \State $X_{\text{hist}}[m+1] \leftarrow x$
            \Statex\hrulefill
            \If{$m \ge n \wedge \twonorm{x - X_{\text{hist}}[m-n]} < \epsilon_{\text{stall}}$} \algcomment{Stalling detected per Definition~\ref{def:stalling}}
                \State $x_{\text{stall}} \leftarrow x$
                \State $N_{stall} \leftarrow \min_{j \in \mathcal{S}} \left(j+\left\lceil k_j / \left(n(c_j - \trans{f_j} x_{\text{stall}})\right)\right\rceil\right)$ \algcomment{Compute using Thm.~\ref{thm:nstall}}
                \State $N_{\text{cycles}} \leftarrow \lceil N_{stall} / n \rceil$
                \For{$j \in \{0, \dots, n-1\}$}
                    \State $k_j \leftarrow \trans{f_j} e_j$ \algcomment{Get scalar from $e_j = k_j f_j$}
                    \State $\Delta k_j \leftarrow \trans{f_j} x_{\text{stall}} - c_j$
                    \State $e_j \leftarrow (k_j + N_{\text{cycles}} \times \Delta k_j) f_j$ \algcomment{Fast-forward dual variable}
                \EndFor
                \State $m \leftarrow m + N_{stall}$ \algcomment{Jump iteration counter}
            \Else
                \State $m \leftarrow m + 1$ \algcomment{Increment iteration counter}
            \EndIf
            \Statex\hrulefill
        \EndWhile
        \State \Return $x$
    \EndFunction
\end{algorithmic}
\end{algorithm}

% \begin{cor}[Convergence of Algorithm~\ref{alg:fastforward}]
% \label{cor:convergence}
% The sequence of iterates $\{x_m\}$ generated by the Stall-Averse Dykstra's Algorithm (Algorithm~\ref{alg:fastforward}) is guaranteed to converge to the optimal solution $x^\star$ asymptotically.
% \end{cor}
% \begin{pf}
% The fast-forwarding step (lines 13--18) is not an approximation. It computes the exact state of the dual variables $\{e_j\}$ that would be attained by the standard Dykstra's algorithm~\eqref{eq:dykstra} after $N_{stall}$ iterations. Since the proposed algorithm achieves the same iterates after escaping the stalling period as would be attained with standard Dykstra, it retains the convergence guarantees of the original Boyle-Dykstra theorem~\cite{DYKSTRA,DYKSTRAPERKINS}. \qed
% \end{pf}
\section{Numerical Example} \label{sec:example}

% We demonstrate the efficacy of our method through a numerical experiment. The results of our experiment can be reproduced from the repository: \href{https://github.com/ClouD-161803/Dykstra-Project.git}{ClouD-161803/Dykstra-Project}.

We demonstrate the efficacy of our method through a numerical experiment. The results of our experiment can be reproduced from the repository~\cite{DykstraProjectRepo}. The example we used is the projection of an initial point $x^\circ$ on the intersection of a box and a line in $\mathbb{R}^2$. The box is centered at the origin and can be represented as $[-1, 1] \times [-1, 1]$, and the line passes through the points $(0,1)$ and $(-2, 2)$. This configuration was first described in detail in~\cite{DYKSTRASTALLING}. The initial point was set to $x^\circ = (-4, 1.4)$ to induce the stalling period, and the ground-truth optimal solution $x^\star$ was computed using the interior-point CQP solver \texttt{quadprog} to serve as a benchmark for error calculations. Half-space activity was monitored by verifying~\eqref{eq:sets_i} for each set at each iteration, and the convergence behaviour monitored via the $\ell_2$ norm of the absolute distance to optimality, $E(x_m) \eqdef \| x_m - x^\star \|_2^2$.
% \begin{align}
% E(x_m) \eqdef \| x_m - x^\star \|_2^2.
% \end{align}
% Half-space activity was monitored by verifying~\eqref{eq:sets_i} for each set at each iteration, and we encoded the status of each half space using an indicator function
% \[
% \mathcal{I}_{\mathcal{H}_i}(x) \eqdef
% \begin{cases}
% 1 & \text{if } x \in \mathcal{H}_i, \\
% 0 & \text{if } x \notin \mathcal{H}_i.
% \end{cases}
% \]

The results are shown in Fig.~\ref{fig:stalling2} for both Dykstra's method and Algorithm~\ref{alg:fastforward}. The iterates $x_m$ are shown in the top plot, the errors $E(x_m)$ in the middle plot and the activity of the half space corresponding to the left side of the box in the bottom plot. Note that the sequence of iterates $x_m$ generated by Dykstra's method and Algorithm~\ref{alg:fastforward} are identical, so they cannot be distinguished in the top plot. However, the error indicates that Dykstra's algorithm is stalling up to cycle 16 ($m\times n$), with the stalling trajectory highlighted in red in the top plot. The half space activity shows that stalling ends when the half space corresponding to the left side of the box becomes inactive.

The middle and bottom plot of Figure~\ref{fig:stalling2} illustrate the application of Algorithm~\ref{alg:fastforward} to the same example, effectively resolving the stalling issue. In Fig.~\ref{fig:stalling2}, stalling is detected at cycle 1, upon which the algorithm is fast-forwarded using $N_{stall}=15$ cycles, where $N_{stall}$ is computed using Thm.~\ref{thm:nstall}. At cycle 2, the iterates of the modified algorithm arrive where the iterates of the original algorithm arrive at iteration 16. The iterates of the modified algorithm then become identical to those of the original algorithm. It is evident that the modified method achieves superior convergence characteristics when compared to standard Dykstra's algorithm. Although this section has showcased a resolution for a simple example in $p = 2$ dimensions with only three active half-spaces, the results of Thm.~\ref{thm:nstall} and Algorithm~\ref{alg:fastforward} will hold in any number of dimensions and for an arbitrary number of half-spaces.

% \begin{figure}[h]
%     \centering
%     \includegraphics[width=.9\textwidth]{Latex/Current Version/Figures/No_stalling_vertical.png}
%     \caption{Modified algorithm with fast-forwarding. (Top): Stalling is detected at iteration 1. The algorithm computes $N_{stall}=15$ using Theorem~\ref{thm:nstall} and ``fast-forwards" the iterates, jumping directly to the point where the standard algorithm would have been at iteration 16. Asymptotic convergence follows from iteration 2 onwards. (Middle): The stalling period (yellow) is eliminated. The algorithm ``converges" (green) at iteration 16. (Bottom): Half-space 0 (black circles) becomes inactive immediately after the fast-forward step.}
%     \label{fig:fastforwarding}
% \end{figure}

% \begin{figure}
%     \centering
%     \begin{subfigure}{\textwidth}
%         \centering
%         \includegraphics[width=0.65\textwidth]{Latex/Current Version/Figures/Stalling_vertical.png}
%         \caption{Stalling situation.}
%         \label{fig:stalling2}
%     \end{subfigure}
%     \begin{subfigure}{\textwidth}
%         \centering
%         \includegraphics[width=0.65\textwidth]{Latex/Current Version/Figures/No_stalling_vertical.png}
%         \caption{Modified algorithm with fast-forwarding.}
%         \label{fig:fastforwarding}
%     \end{subfigure}
%     \caption{Dykstra's algorithm applied to a line-box example. The first column illustrates the trajectories of the iterates and their distance to the Euclidean projection computed using an interior point method. The second column shows the half-space activity, where half-space 0 corresponds to the left side of the box, half-space 1 to the top side of the box, and half-space 0 to the line (down-facing normal).}
%     \label{fig:stallingfixed}
% \end{figure}

\definecolor{legreddraw}{HTML}{E41A1C}
\definecolor{legbluedraw}{HTML}{377EB8}
\definecolor{leggreendraw}{HTML}{4DAF4A}
\definecolor{legredfill}{HTML}{B61516}
\definecolor{legbluefill}{HTML}{2C6593}
\definecolor{leggreenfill}{HTML}{3E8C3B}

\begin{figure}[h]
    \centering
    % \includegraphics[width=.9\textwidth]{Latex/Current Version/Figures/Stalling_vertical.png}
\begin{tikzpicture}
  % Top axis
  \begin{axis}[
    at={(0,0)},
    anchor=north west,
    width=0.8\linewidth,
    height=4cm,
    xmin=-4, xmax=1,
    ymin=0.5, ymax=1.5,
    xlabel={$x$ coordinate},
    ylabel={$y$ coordinate},
    grid=both,legend style={fill=white,fill opacity=0.4,text opacity=1,draw=black,cells={anchor=west}},
    legend pos=south west, legend columns=3, transpose legend,
  ]
    \draw[fill=legbluefill, fill opacity=0.2, draw=legbluedraw]
        (axis cs: -1,0.5) rectangle (axis cs: 1,1);
    \addlegendimage{area legend,fill=legbluefill,fill opacity=0.2,draw=legbluedraw}
    \addlegendentry{Box}

    \draw[very thick, legbluedraw] (axis cs:-2,2) -- (axis cs:1,0.5);
    \addlegendimage{line legend,very thick, draw=legbluedraw}
    \addlegendentry{Line}
    
    \draw[line width=2, legredfill!30, line join=round] (axis cs:-1,1.4) -- (axis cs:-0.8,1.4) -- (axis cs:-1,1) -- (axis cs:-1,1.4);
    \addlegendimage{line legend,line width=2, draw=legredfill!30}
    \addlegendentry{Stalling}
        
    \addplot[
        black,
        thin, densely dotted, line cap= round,
        mark=*, mark size = 0.75, line join=round,
    ] table[x=x, y=y, col sep=comma] {data_iterates.csv};
    \addlegendentry{Iterates $x_m$}

    \addplot[
    orange,
    only marks,
    mark=star,
    mark size=2.5,
    ] coordinates {(0, 1)};
    \addlegendentry{Projection $x^\star$}

  \end{axis}

  % Middle axis
  \begin{axis}[
    at={(0,-3.75cm)},
    anchor=north west,
    width=0.8\linewidth,
    height=4cm,
    xmin=0, xmax=30,
    ymode=log,
    ylabel={Squared error},
    xlabel={Cycle},
    grid=both,
    legend pos=south west,
    legend style={fill=white,fill opacity=0.4,text opacity=1,draw=black,cells={anchor=west}},
    legend entries={Original, Algorithm~\ref{alg:fastforward}},
  ]
    \addplot[
        black,
        thick,
        mark=*, mark size= 1,
    ] table[x expr=\thisrow{iteration} + 1, y=error, col sep=comma] {data_dykstra_original.csv};
    \addplot[
        leggreendraw,
        thick,
        mark=x, mark size= 2,
    ] table[x expr=\thisrow{iteration} + 1, y=error, col sep=comma] {data_dykstra_modified.csv};
  \end{axis}

  % Bottom axis
  \begin{axis}[
    at={(0,-7.5cm)},
    anchor=north west,
    width=0.8\linewidth,
    height=4cm,
    xmin=0, xmax=30,
    ylabel={Halfspace activity},
    xlabel={Cycle},ytick={0,1},
    grid=both,
    legend style={fill=white,fill opacity=0.4,text opacity=1,draw=black,cells={anchor=west}},
    legend entries={Original, Algorithm~\ref{alg:fastforward}},
  ]
    \addplot[
        black,
        thick,
        mark=*, mark size= 1,
    ] table[x expr=\thisrow{iteration} + 1, y=halfspace, col sep=comma] {data_dykstra_original.csv};
    \addplot[
        leggreendraw,
        thick,
        mark=x, mark size= 2,
    ] table[x expr=\thisrow{iteration} + 1, y=halfspace, col sep=comma] {data_dykstra_modified.csv};
  \end{axis}
\end{tikzpicture}
    \caption{Dykstra's projection algorithm and Algorithm~\ref{alg:fastforward} applied to the line-box example from~\cite{BAUSCHKESWISS}. Top: The iterates $x_m$ of the algorithms with stalling of Dykstra's method highlighted in red. Middle: Squared errors $E(x_m)$. (Bottom): Activity of the vertical half space (left side of the box). The stalling period ends when this half space becomes inactive.}
    \label{fig:stalling2}
\end{figure}
\section{Conclusion}
\label{sec:conclusion}

Dykstra's algorithm is a method for solving the Euclidean BAP on the intersection of convex sets, but the stalling phenomenon has significantly hindered its practical utility~\cite{BAUSCHKESWISS}. For time-critical applications, such as optimisation in real-time control contexts, an algorithm with an unpredictable and arbitrarily long execution time is not practical. In this paper, we have proposed a solution to the stalling problem for the polyhedral case. We first formalised the stalling period (Definition~\ref{def:stalling}) and then derived its exact length in closed-form (Theorem~\ref{thm:nstall}), rendering the stalling period computable \textit{a priori}, once stalling is detected. Based on this result, we proposed the Algorithm~\ref{alg:fastforward}. This modified algorithm detects stalling and uses the result from Thm.~\ref{thm:nstall} to ``fast-forward'' the algorithm's auxiliary variables past the entire stalling period in a single, computationally inexpensive step. As demonstrated by our numerical results in Section~\ref{sec:example}, the modification eliminates the stalling behaviour (Fig.~\ref{fig:stalling2}). Since the proposed algorithm produces the same iterates after escaping the stalling period as would be produced with standard Dykstra, the Boyle-Dykstra theorem is retained and the iterates are guaranteed to converge to the optimal solution $x^\star$ asymptotically (Corollary~\ref{cor:convergence}).
Our method ensures the iterates converge faster, which enables practitioners to implement Dykstra's algorithm.

We identify several directions for future research that are worth exploring. First, extending the analysis from polyhedral sets to broader classes of convex sets would enhance the generality of the results. Second, applying our method to a reformulated linear program, where the optimality conditions are decomposed into a convex and a conic set, could provide an interesting benchmark for evaluating the proposed approach. Finally, although we have improved Dykstra's overall convergence properties by eliminating stalling, we have not been able to show that our modified version yields iterates $x_m$ that monotonically approach the projection $x^\star$. According to the analysis in~\cite{DYKSTRAPOLY}, the algorithm successively discards sets $\mathcal{H_i}$ for which $x^\star\in\text{int}\mathcal{H_i}$, and during these phases, the algorithm can produce iterates that temporarily diverge from $x^\star$. Finding a modification that fast-forwards these phases too would allow one to bound the error $\anynorm{x_m-x^\star}$ \textit{a priori}, and guarantee that the output of Algorithm~\ref{alg:fastforward} is closer to $x^\star$ when executed for a fixed number of iterations.



\begin{ack}
This work was supported by Keble College using the Keble College Small Research Grant \#KSRG118.
\end{ack}

\printbibliography

\end{document}